\documentclass[11pt]{article}
\usepackage[margin=1in]{geometry}
\usepackage{amsmath,amssymb,amsthm,mathtools,mathrsfs}
\usepackage{graphicx}
\usepackage[colorlinks=true,linkcolor=blue,citecolor=blue,urlcolor=blue]{hyperref}
\usepackage{microtype}

\newtheorem{theorem}{Theorem}
\newtheorem{lemma}{Lemma}
\newtheorem{remark}{Remark}
\newcommand{\dist}{\operatorname{dist}}
\newcommand{\M}{\mathcal M}
\newcommand{\N}{\mathscr N}
\makeatletter
\renewcommand{\@seccntformat}[1]{\csname the#1\endcsname\hspace{0.8em}}
\makeatother

\title{The stated system-of-trivial-centralizers question\\
is universally affirmative}
\author{Run \texttt{fa\_banach\_001}, lane 18
  (\texttt{agent\_lane\_18})}
\date{13 August 2026}

\begin{document}
\maketitle

\begin{abstract}
The first concluding question of Cabello S\'anchez's paper asks whether each
$\M$-centralizer $\Omega:L^p_0\to L^q$, $0<q<p<\infty$, admits a system of
trivial centralizers whose distances from the module morphisms have supremum
equal to the corresponding distance of $\Omega$.  As written, the question
places no relation between $\Omega$ and the proposed system.  We prove that
this omission makes the answer universally affirmative for arbitrary
quasinormed modules and all exponents.  If $\Omega$ has finite distance, use
$\Omega$ itself.  If its distance is infinite, scale one bounded nonlinear
homogeneous map having positive distance from the morphisms.  Such maps are
trivial centralizers.  This resolves the literal display but does not produce
the intended analogue of the paper's finite-rank localization lemma.
\end{abstract}

\section{Exact source question}

The source asks the following in its concluding remarks:
\begin{quote}
Most results in Sections 3 and 4 would generalize to noncommutative $L^p$
spaces associated to arbitrary von Neumann algebras as long as one could find
a good substitute for Lemma 3.  More precisely, we ask if for every
$\M$-centralizer $\Omega:L^p_0\to L^q$ with $0<q<p<\infty$ there is a system
of trivial centralizers $\Omega_i$ such that
\[
 \dist(\Omega,\mathscr{M}_{\M}(L^p_0,L^q))
 =\sup_i\dist(\Omega_i,\mathscr{M}_{\M}(L^p_0,L^q)).
\]
Here, $L^p_0=\{af^{1/p}:a\in\M\}$, where $f$ is a normal, faithful state on
$\M$.
\end{quote}
The exact source page is reproduced below.
\begin{center}
\includegraphics[width=0.98\textwidth]{figures/open_problem_crop.png}
\end{center}

The preceding Lemma 3 uses the concrete localizations
$\Phi_e(x)=\Phi(xe)$ for finite-rank projections $e$.  The displayed question,
however, does not say that $\Omega_i$ must be a localization of $\Omega$, nor
does it impose any compatibility, convergence, or uniform centralizer bound.
The proof below takes the quantifiers exactly as displayed.

\section{Definitions and proof intuition}

For homogeneous maps $F:X\to Y$ between quasinormed spaces, set
\[
 \|F\|_{\mathrm h}=\sup_{x\ne0}\frac{\|F(x)\|_Y}{\|x\|_X}
 \quad\text{and}\quad
 \dist(F,\N)=\inf_{T\in\N}\|F-T\|_{\mathrm h},
\]
with values in $[0,\infty]$.  Here $\N$ is the vector space of algebraic
module morphisms.  A centralizer is trivial precisely when its distance from
$\N$ is finite.

The key observation is numerical.  The question asks for a family of trivial
maps realizing the same extended number $D=\dist(\Omega,\N)$, but it does not
require those maps to contain any information about $\Omega$.  Finite $D$ is
realized by $\Omega$ itself.  Infinite $D$ is realized by scalar multiples of
one unrelated trivial map whose distance is positive.

\section{A trivial centralizer at positive distance}

We use only the standard assumptions that the module actions are jointly
continuous and that the quasi-norms are homogeneous and Hausdorff.

\begin{lemma}\label{positive-map}
Let $X$ and $Y$ be nonzero complex quasinormed left modules over a Banach
algebra $A$, with $\dim_{\mathbb C}X>1$.  Let
$\N=\mathscr{M}_A(X,Y)$.  There is a bounded homogeneous map
$B:X\to Y$ such that:
\begin{enumerate}
\item $B$ is quasilinear and an $A$-centralizer;
\item $B$ is a trivial centralizer;
\item $0<\dist(B,\N)<\infty$.
\end{enumerate}
\end{lemma}

\begin{proof}
Choose a unit representative $s_L$ on each complex one-dimensional subspace
$L\subset X$.  Fix two linearly independent unit vectors $u,v$, and choose
$s_{\mathbb Cu}=u$.  Every nonzero $x\in X$ has a unique representation
$x=\lambda s_L$.  Define
\[
 b(0)=0,\qquad
 b(\lambda s_L)=
 \begin{cases}
   \lambda,&L=\mathbb Cu,\\
   0,&L\ne\mathbb Cu.
 \end{cases}
\]
Then $b$ is complex homogeneous and $|b(x)|\leq\|x\|_X$.  It is not additive:
$b(u)=1$, $b(v)=0$, whereas $b(u+v)=0$.  Fix a unit vector $y_0\in Y$ and put
$B(x)=b(x)y_0$.  Thus $B$ is bounded, homogeneous, and not additive.

Let $K_X,K_Y$ be quasi-triangle constants and let $C_X,C_Y$ bound the two
module actions.  Bounded homogeneity gives
\begin{align*}
 \|B(x+z)-B(x)-B(z)\|_Y
 &\leq C\bigl(\|x+z\|_X+\|x\|_X+\|z\|_X\bigr)\\
 &\leq C'\bigl(\|x\|_X+\|z\|_X\bigr),
\end{align*}
so $B$ is quasilinear.  Likewise,
\begin{align*}
 \|B(ax)-aB(x)\|_Y
 &\leq C\bigl(\|B(ax)\|_Y+\|aB(x)\|_Y\bigr)\\
 &\leq C'\|a\|_A\|x\|_X,
\end{align*}
so $B$ is a left centralizer.  Since $B$ is bounded, the zero morphism
witnesses triviality and $\dist(B,\N)\leq\|B\|_{\mathrm h}<\infty$.

It remains to exclude distance zero.  If $\dist(B,\N)=0$, for every $n$ one
can choose $T_n\in\N$ satisfying
\[
 \|(B-T_n)(x)\|_Y\leq n^{-1}\|x\|_X\qquad(x\in X).
\]
Since $T_n$ is additive, the quasi-triangle inequality yields, for fixed
$x,z\in X$,
\begin{align*}
 \|B(x+z)-B(x)-B(z)\|_Y
 &\leq \frac{C}{n}
  \bigl(\|x+z\|_X+\|x\|_X+\|z\|_X\bigr).
\end{align*}
Letting $n\to\infty$ forces $B(x+z)=B(x)+B(z)$, contradicting the
construction.  Hence $\dist(B,\N)>0$.
\end{proof}

No local convexity is used; in particular, the lemma applies when the target
is $L^q$ with $0<q<1$.

\section{Universal answer to the displayed question}

\begin{theorem}\label{main}
Let $X,Y$ be complex quasinormed modules and let
$\N=\mathscr{M}_A(X,Y)$.  For every centralizer $\Omega:X\to Y$, there exists
a system of trivial centralizers $(\Omega_i)$ such that
\[
 \dist(\Omega,\N)=\sup_i\dist(\Omega_i,\N).
\]
Consequently, the first concluding question of arXiv:1910.08774 has an
affirmative answer exactly as written, for every $p,q>0$ and every von
Neumann algebra $\M$.
\end{theorem}

\begin{proof}
Put $D=\dist(\Omega,\N)$.

If $D<\infty$, then by definition there is a morphism at finite distance from
$\Omega$, so $\Omega$ is a trivial centralizer.  Take the one-element system
$\Omega_1=\Omega$.  Its supremal distance is exactly $D$.

Suppose $D=\infty$.  Then $\Omega$ cannot be bounded, because the zero
morphism would otherwise give finite distance.  In particular $X$ cannot be
one-dimensional: every complex homogeneous map on a one-dimensional
quasinormed space is bounded.  Apply Lemma~\ref{positive-map} and put
\[
 d=\dist(B,\N),\qquad \Theta=d^{-1}B.
\]
The class $\N$ is a complex vector space, hence distance is homogeneous:
\[
 \dist(\lambda B,\N)=|\lambda|\dist(B,\N).
\]
Therefore every $\Omega_n=n\Theta$ is a trivial centralizer and
\[
 \dist(\Omega_n,\N)=n.
\]
It follows that
$\sup_n\dist(\Omega_n,\N)=\infty=D$, as required.
\end{proof}

\section{Why this is not the intended Lemma 3 substitute}

The theorem is deliberately overgeneral: it uses neither $q<p$ nor any
property of von Neumann algebras.  This shows that the displayed formulation
does not encode the local-to-global mechanism described immediately before
it.

In the Schatten setting, the local maps are prescribed by
$\Phi_e(x)=\Phi(xe)$, where $e$ ranges over finite-rank projections.  Right
multiplication commutes with the left action, the compressed operator lies in
the finite-rank core, and $\|x\|_\infty\leq\|x\|_p$ makes every $\Phi_e$
trivial.  An ultrafilter and WOT lower semicontinuity then recover the global
distance.

For $L^p_0=\M f^{1/p}$, an arbitrary right projection need not preserve this
core.  Modular-analytic right multipliers can preserve it, but the elementary
centralizer estimate controls $\|a\|_\infty$ rather than
$\|af^{1/p}\|_p$, so the crucial Schatten bound still does not follow.  The
argument above says nothing about this obstruction.

A meaningful repaired question must specify, at minimum:
\begin{enumerate}
\item how every $\Omega_i$ is obtained from the given $\Omega$;
\item which local pieces of $L^p_0$ the index set exhausts;
\item the compatibility needed to patch the approximating morphisms; and
\item any required uniform bound on centralizer constants.
\end{enumerate}
With such conditions imposed, Theorem~\ref{main} no longer applies, and the
substantive arbitrary-von-Neumann-algebra problem remains unresolved here.

\section{Verification and novelty audit}

The proof reduces to four exact checks:
\begin{enumerate}
\item the chosen-line map $b$ is complex homogeneous, bounded, and nonadditive;
\item bounded homogeneous maps are quasilinear centralizers and are trivial
via zero;
\item zero distance from the morphisms would force additivity;
\item distance from a vector subspace is homogeneous under scalar dilation.
\end{enumerate}
Each check is written out above and remains valid for quasi-norms.

Exact-phrase searches, arXiv searches for the displayed $L^p_0$ question, the
run's result indexes, and later papers citing the source did not locate a
published correction or answer.  This is evidence rather than a definitive
citation audit.  The mathematical claim is a full answer only to the literal
display.  Before dissemination, a specialist should verify that the author
did not intend a convention in which the word ``system'' silently includes
compression/localization data.  If no such convention exists, the useful
conclusion is that the question needs an additional hypothesis, not that the
intended extension theory has been solved.

\begin{thebibliography}{9}
\bibitem{source}
F.~Cabello S\'anchez,
\emph{Nonlinear centralizers in homology II. The Schatten classes},
Rev. Mat. Iberoam. 37 (2021), no.~6, 2309--2346;
arXiv:1910.08774.

\bibitem{twisting}
F.~Cabello S\'anchez, J.~M.~F. Castillo, S.~Goldstein, and
J.~Su\'arez de la Fuente,
\emph{Twisting non-commutative $L_p$ spaces},
Adv. Math. 294 (2016), 454--488; arXiv:1407.0969.
\end{thebibliography}

\end{document}
